% ------------------------------------------------------------------
% Exemplo de introdução gerada por textos dummys a partir do
% lipsum
%-------------------------------------------------------------------


\chapter{Introdução}
\label{cap:intro} % faço a referência na bibliografia
Ao final do século XIX, em meio à crescente atividade econômica impulsionada pela Revolução Industrial, William Stanley Jevons publicou o livro \textit{The Coal Question} \citeyear{Jevons1866}, no qual aborda a crescente preocupação da época quanto ao exaurimento das minas de carvão, principal combustível dos crescentemente importantes motores a vapor. O que passou a chamar-se ``paradoxo de Jevons'' é um fenômeno por ele observado em sua publicação: uma maior eficiência do emprego de carvão induz um \textit{aumento} de seu consumo, oposto à intuição cotidiana que ela induziria uma \textit{diminuição} de seu consumo (pois menos energia seria necessária para alcançar um mesmo nível de produto). 
%Em uma famosa passagem, Jevons afirma que ``É uma confusão de ideias supor que o uso econômico de combustível [carvão] seja equivalente à diminuição do [seu] consumo. É o contrário que é verdade.'' \cite[p. 75]{Jevons1866}.

Jevons trata, portanto, do aumento \textit{absoluto} do consumo de carvão enquanto fonte energética, não obstante a redução \textit{relativa} da necessidade de seu consumo para um mesmo nível prévio de produção --- isso pode ser constatado, por exemplo, pelo emprego paulatino de motores com maior eficiência termodinâmica, os quais foram usados em escalas crescentes de produção (e, portanto, de consumo de matéria-prima e de energia). Ou seja, embora tais ganhos de eficiência \textit{permitam} um consumo menor de carvão para um mesmo nível de produção prévio, tende a ocorrer um aumento do nível de produção \textit{efetivo}. De certa forma, a possibilidade de poupar recursos dá espaço para seu consumo alhures, conduzindo a uma maior atividade econômica e, portanto, maior consumo efetivo de carvão do que previamente.

Houve também desenvolvimentos posteriores no que tange à caracterização deste ``paradoxo'' que foram levados a cabo por economistas neoclássicos, quando das crises do petróleo na década de 1970, cerca de um século após a análise de Jevons. Os trabalhos basilares neste tema são de \textcite{Khazzoom1980, Khazzoom1987}, \textcite{Brookes1990} e \textcite{Saunders1992}, este último sendo o responsável por cunhar a expressão ``postulado de Khazzoom-Brookes'': 
\begin{quote}
    \textit{Postulado de Khazzoom-Brookes}: Com um preço real da energia fixo, os ganhos em eficiência energética aumentarão o consumo de energia acima do que seria sem esses ganhos. \parencite[p. 135]{Saunders1992}
\end{quote}

O enunciado sobre restrição de ``preços reais da energia fixos'' é relevante, pois permite (em tese) isolar ganhos de eficiência \textit{de facto} de meras mudanças de preços que possam aparecer como ``ganhos/perdas de eficiência''. Brookes, a título de exemplo, menciona que o choque do petróleo em 1973 ``levou 4\% do parque industrial americano à falência'' --- justamente os setores mais dependentes de energia e, portanto, os mais afetados pelos aumentos drásticos do preço do petróleo ---, o que levou ao que poderia ser interpretado como um ``aumento de eficiência energética'' do parque industrial resultante --- justamente pelo cessar das atividades destes setores mais intensivos em energia ---, ``sem que ninguém necessariamente tivesse adotado processos e práticas mais eficientes em termos energéticos'' \parencite[p. 323]{Brookes1990}. 

Na prática, porém, tais movimentos --- variações de eficiência material e variações de preços envolvidos --- são indissociáveis, cristalizando-se no que a literatura de Economia da Energia passou a chamar de ``efeito \textit{rebound}''\footnote{Alguns o traduzem como ``efeito rebote'' ou ``efeito bumerangue''. Priorizo o nome no original, em inglês, para que não dê margem para equívocos.}, que trata sobre a variação do consumo energético após algum ``choque tecnológico'' que permita maior (ou demande menor) poupança de energia. O efeito \textit{rebound} é uma espécie de ``efeito preço'' microeconômico\footnote{A título de ilustração, supõe-se um caso em que há um \textit{aumento} na eficiência energética para ilustração da decomposição do ``efeito preço'' do consumo energético. \textit{Mutatis mutandis} para o caso de sua diminuição.}: é a composição de um ``efeito substituição'', pelo qual permite-se uma diminuição no consumo do fator energético de produção, e um ``efeito produto'', que expande a fronteira de possibilidade de produção e, portanto, sua demanda/consumo \cite[p. 174]{Saunders2009}.

No que tange à observação efetiva deste fenômeno, Sorrell destaca o escopo no qual julga-se adequado avaliar empiricamente o efeito \textit{rebound}: 

\begin{quote}
    Para o postulado de K-B [Khazzoom-Brookes], o limite relevante do sistema é normalmente considerado a economia nacional. No entanto, \textit{melhorias na eficiência energética também podem afetar os padrões comerciais e os preços internacionais da energia, alterando assim o consumo de energia em outros países}. [...] Para capturar toda a gama de efeitos \textit{rebound}, o limite do sistema para a variável independente (eficiência energética) deve ser relativamente estreito, enquanto o limite do sistema para a variável dependente (consumo de energia) deve ser o mais amplo possível. \cite[p. 15, grifo meu]{Sorrell2007}.
\end{quote}

O postulado de Khazzoom-Brookes, como o paradoxo de Jevons, trata-se de quando o efeito \textit{rebound} é grande o suficiente para que a poupança energética \textit{esperada}, tornada possível pelo choque tecnológico, é totalmente consumida. Dito de outra forma, a poupança \textit{efetiva} ou é nula --- i.e. não verifica-se diferença no consumo energético ---, ou mesmo ``negativa'', consumindo-se \textit{mais} energia do que antes do aumento de eficiência. Neste último caso, diz-se que houve um ``efeito \textit{backfire}''\footnote{Ocasionalmente traduzido como ``efeito tiro-pela-culatra''. Vide nota acima, mantenho o termo no original por clareza.}. O postulado de Khazzoom-Brookes, portanto, prevê que ganhos de eficiência energética induzem (ou \textit{tendem a induzir}) um efeito \textit{backfire} do consumo energético, ou seja, um consumo \textit{efetivo} maior do que o consumo mínimo \textit{esperado}.



%----------------------------------------------------------------------
\chapter{Justificativa}
\label{cap:justificativa}
\section{Motivação}
Análises empíricas sobre o efeito \textit{rebound/backfire} acabam tendo certas limitações teóricas. Pode-se perceber, por exemplo, que há uma conflação entre o consumo energético \textit{produtivo/industrial} e o consumo \textit{improdutivo/domiciliar} em suas descrições neoclássicas. Já aqui, em particular, há uma diferença fundamental com o fenômeno que Jevons observou ao final do século XIX, pois ele escreve explicitamente que ``a economia [de carvão] obtida [por usos domésticos] não ultrapassaria cerca de um terço do consumo total de carvão''. A famosa citação de Jevons, em contrapartida, trata-se do parágrafo seguinte ao supracitado, embora raramente seja reproduzida \textit{in totum}, vide abaixo: 
\begin{quote}
    Mas a economia do carvão \textit{nas manufaturas} é uma questão diferente. É uma completa confusão de ideias supor que o uso econômico de combustível seja equivalente a um consumo reduzido. O contrário é que é verdade. \cite[p. 75]{Jevons1866}
\end{quote}

A primeira frase, que restringe esta afirmação \textit{ao consumo manufatureiro} de carvão, costuma ser omitida em citações pós-século XX, embora fosse relevante ao consumo da época de Jevons. Ainda assim, por mais que o consumo domiciliar de energia tenha crescido sobremaneira desde o final do século XIX --- consistindo essencialmente em maior adoção residencial de energia elétrica  ---, mesmo assim ele tende a ser menor do que o consumo industrial de energia\footnote{Cf. \textcite{USEnergyInformation}, por exemplo.}. 

Não obstante seus usos serem qualitativamente distintos --- o uso domiciliar é (ou tende a ser) improdutivo, enquanto seu uso industrial é empregado para fins produtivos e geração de lucros ---, seus usos tendem a ser misturados em análises do efeito \textit{rebound}. Ofuscam-se, portanto, as distinções qualitativas destes usos, podendo ocorrer situações em que o caso que a diminuição de energia elétrica pelo setor domiciliar seja interpretada como poupança do setor privado, ou em que o aumento dos gastos deste último sejam passíveis de ``se tornar culpa'' daquele primeiro. 


%\begin{quote}
%    É muito comum argumentar que a escassez de carvão será suprida por novas formas de utilizá-lo de maneira eficiente e econômica. A quantidade de trabalho útil obtida a partir do carvão pode ser multiplicada, enquanto o consumo permanece estável ou diminui. Supõe-se, assim, que teremos os meios para neutralizar completamente os males da escassez e do alto custo desse combustível. [...]

%    \textit{Não me refiro aqui ao consumo doméstico de carvão}. Este, sem dúvida, pode ser reduzido sem maiores prejuízos além da diminuição do conforto em nossos lares e de uma certa alteração em nossos hábitos nacionais consolidados. O carvão assim economizado seria, em sua maior parte, armazenado para uso das gerações futuras. Mas mesmo que nossa população pudesse ser induzida a se abster do prazer de uma boa lareira, a economia obtida não ultrapassaria cerca de um terço do consumo total de carvão [...] Dos outros dois terços, quase um terço é utilizado em nossas fábricas de ferro; e o restante em nossas fábricas, fornos e oficinas mecânicas em geral. \parencite[p. 75, grifo meu]{Jevons1866}
%\end{quote}


\section{Relevância}
A discussão de sustentabilidade sobre transição energética frequentemente cai num discurso de ``otimismo tecnológico'', em que presume-se que problemas estruturais de hoje podem ser resolvido através de avanços tecnológicos; um exemplo em voga hoje em dia é a questão da transição da matriz energética rumo a energias sustentáveis. Contudo, mesmo a economia ortodoxa, que tende a ser ``apologética'' às estruturas capitalistas\footnote{Uma discussão detalhada sobre tal caráter da economia ortodoxa (e mesmo heterodoxa) pode ser encontrado em \textcite{Medeiros2013}.}, descobre a contradição deste discurso, como apontado (em particular) por \textcite{Jevons1866} e \textcite{Saunders1992}: a mera busca pelo aumento da eficiência energética não só pode não auxiliar na diminuição de seu consumo, como pode até mesmo \textit{aumentá-lo}.

% O objetivo deste trabalho é fundamentar tal visão pela teoria marxista do valor-trabalho, a fim de elucidar a necessidade lógica deste fenômeno --- aumentos em eficiência energética são (tendencialmente) seguidos por aumentos em consumo energético, em um nível macroeconômico --- através da lógica imanente do capital. Ou seja, este trabalho busca desfazer não só seu caráter aparentemente paradoxal, como também a ``naturalidade'' que pode vir a coadunar-se com este ``fato social''. 

% A descrição do efeito \textit{rebound/backfire} pela economia ortodoxa comete não só a conflação entre consumos produtivo e improdutivo, mas, e talvez mais pertinentemente, a conflação entre o consumo \textit{individual} e o consumo \textit{social} --- ou, em termos mais usuais na Economia, entre fenômenos microscópicos e macroscópicos. Embora o efeito \textit{rebound} seja definido como um fenômeno emergente, advindo das várias interações dentro do sistema complexo que é a economia como um todo, não é incomum que autores que o abordam utilizem-se de analogias que se referem a comportamentos \textit{individuais}. Por exemplo, \textcite{Berkhout2000}, na seção de definição do efeito \textit{rebound}, empregam a analogia 
%\begin{quote}
%    Se um carro, por exemplo, consegue percorrer mais quilômetros com um litro de gasolina, o custo do combustível por quilômetro diminui, assim como o custo total por quilômetro. \parencite[p. 426]{Berkhout2000}
%\end{quote}

%Ou seja, tendem a incorrer na falácia da composição: ``[acreditam] que a redução do consumo de energia no nível microeconômico de aplicações individuais levará automaticamente a reduções proporcionais no nível macroeconômico'' \parencite{Brookes1990}. 

Ainda no ínterim do ``otimismo tecnológico'', frequentemente adere-se à visão de que o aumento da eficiência tecnológica, ao necessitar de menos recursos para seu funcionamento, teoricamente gerará menos impactos sobre o meio ambiente, em certos casos até mesmo alegando a possibilidade de ``desmaterialização'' da economia \parencite{SaBarreto2016}. Porém, como discutido até aqui, está claro que tais avanços tecnológicos, embora sejam importantes para a economia e o bem-estar da população, não são suficientes \textit{per se} no que tange à mitigação da crise climática, porquanto a questão do consumo energético, inserido no modo de produção capitalista, tem uma tendência imanente a ser fomentado conforme aumenta-se a eficiência (principalmente produtiva) de seu emprego; tendem, portanto, a fomentar a produção capitalista, a extração de recursos e suas consequências socioambientais. 





%----------------------------------------------------------------------
\chapter{Objetivos}
\label{cap:objetivos}
\section{Objetivos gerais}
Este trabalho busca compreender o postulado de Khazzoom-Brookes --- ou melhor, o fenômeno socioeconômico que ele busca descrever --- através da teoria do valor-trabalho, exposta por Marx n'\textit{O Capital} \citeyear{Marx2017a, Marx2014, Marx2017}, valendo-se também de trabalhos posteriores relacionados, como em \textcite{SaBarreto2018}. O objetivo é não só entender o porquê deste fenômeno, como também entender a \textit{necessidade} de sua ocorrência, devida à lógica imanente do capital. 

Tal lógica é observável, por exemplo, no ímpeto do capital de aumentar não só a \textit{eficiência} no consumo produtivo de meios de produção --- requerer \textit{menos} destes para um nível de produção dado ---, mas também a \textit{produtividade} da produção --- em um dado período de tempo, produzir mais mercadorias, e, portanto, requerendo \textit{mais} meios de produção \parencite{SaBarreto2021a}. Este é o alicerce sobre o qual assenta-se o fenômeno socioeconômico denominado ``postulado de Khazzoom-Brookes''. 




\section{Objetivos específicos}
Enunciá-lo desta forma, em particular, permite compreender como este fenômeno socioeconômico propaga-se socialmente não só devido a ``hábitos perdulários'' (em casos de aumento de consumo), mas --- e principalmente --- devido ao caráter desmedido da autovalorização do capital. Tendemos a pensar ``o capital'', contudo, como restrito à esfera da produção, relegando a esfera do consumo ao ``consumidor final'', ao ``mercado''. Um dos objetivos de traduzir essa discussão em termos marxistas é de recuperar a relação reflexiva entre produção, distribuição e consumo, analisando-as como momentos da totalidade do movimento do capital \parencite{Marx2021}. Recupera-se, dessa forma, a interconexão entre tais momentos, evitando-se cair na falácia da composição (``microcomportamentos ditam macrocomportamentos'') ou na lei de Say (``oferta dita demanda'', ``produção dita consumo'').

O objetivo principal de tal análise marxista, sob o prisma da teoria do valor-trabalho, é da descrição do movimento imanente do capital, destacando como, em particular, a questão energética assume a caracterização apreendida pela economia ortodoxa como ``paradoxal''. Ou seja, busca não só descrever o vir-a-ser deste fenômeno, como também o porquê da aparência que este assume aos agentes econômicos e, em particular, aos economistas. \parencite{Marx2017,Marx2014,Marx2017a}. 


%----------------------------------------------------------------------
\chapter{Metodologia}
\label{cap:metodologia}
%- Trabalho bibliográfico
%- Principais referências abordadas
%- Principais objetos da crítica e como ela será abordada

%Visão pela economia neoclássica: \textcite{Saunders1992}, \textcite{Brookes1990,Brookes1993,Brookes2000}, \textcite{Sorrell2007,Sorrell2009}, \textcite{Berkhout2000}.

%Visão pela economia política da energia: \textcite{Malm2013}, \textcite{Pasquinelli2022,Pasquinelli2023}

%Visão marxista da teoria do valor trabalho: 

O projeto consiste em um trabalho bibliográfico teórico de reformulação do postulado de Khazzoom-Brookes sob a teoria marxista do valor-trabalho. Busca-se, para tanto, 1) uma descrição materialista de o que é energia (no que tange ao capital), 2) como e quando ela se torna passível de ser manuseada em atividades econômicas --- i.e. como (e quando) ela torna-se mercadoria ---,  e 3) como este ``paradoxo'' é, de fato, uma consequência lógica do movimento imanente do capital. Ou seja, o objetivo é não de meramente ``retraduzir'' as categorias da economia neoclássica quanto a este fenômeno em termos marxistas, mas de desvelar seus aspectos mais essenciais e determinantes e, dessa forma, desvendar por que ele aparece da forma que aparece. 

Para isso, o trabalho dividir-se-á em uma introdução, 4 capítulos, e uma seção de considerações finais.

A introdução proverá uma contextualização da problemática abordada por \textcite{Jevons1866} e, posteriormente, principalmente por \textcite{Saunders1992} e \textcite{Brookes1990,Brookes1993,Brookes2000}, assim como apresentará a problemática que motiva a abordagem marxista deste fenômeno.

O primeiro capítulo fará uma breve descrição sobre o surgimento histórico da categoria ``energia'' no contexto capitalista, em decorrência dos desenvolvimentos advindos da Revolução Industrial e da paulatina substituibilidade de trabalhadores humanos por máquinas em processos industriais de produção, cf. \textcite{Pasquinelli2022,Pasquinelli2023} e \textcite{Malm2013}. A importância deste capítulo é de estabelecer uma base materialista sobre o conceito \textit{energia}, melhor delimitando-o para fins de estudo. 

O segundo capítulo discorrerá sobre o caráter do que será denominado como ``mercadoria-energia'': analogamente a como o usufruto da mercadoria \textit{força de trabalho} trata-se de um trabalho útil particular\footnote{``No mercado, o que se contrapõe diretamente ao possuidor de dinheiro [i.e. o capitalista] não é, na realidade, o trabalho, mas o trabalhador. O que este último vende é sua força de trabalho. \textit{Mal seu trabalho tem início e a força de trabalho já deixou de lhe pertencer}, não podendo mais, portanto, ser vendida por ele. O trabalho é a substância e a medida imanente dos valores, mas ele mesmo não tem valor nenhum.'' \cite[p. 607, grifo meu]{Marx2017}}, o usufruto da mercadoria-energia trata-se de geração de energia \textit{útil}\footnote{Oposto a, por exemplo, energia \textit{primária}.}. A relevância deste capítulo, assim como a discussão do ``valor do trabalho'' em \textcite{Marx2017}, busca desmistificar a noção de ``valor/preço da energia'' que acaba por ofuscar toda a discussão ortodoxa sobre o postulado de Khazzoom-Brookes.

O terceiro e quarto capítulos apresentarão a discussão que Marx faz n'\textit{O Capital} sobre a teoria do valor-trabalho, destacando sua pertinência no que tange à questão da mercadoria-energia, elucidando não só a substância do postulado de Khazzoom-Brookes, como sua necessidade lógica em razão do movimento imanente do capital \cite{Marx2017,Marx2014,Marx2017a,SaBarreto2016,SaBarreto2018,SaBarreto2022}.

A seção de considerações finais fará uma síntese da discussão do postulado de Khazzoom-Brookes à luz da teoria marxista, evidenciando como este fenômeno aparece, e por que aparece desta forma, para os agentes econômicos em geral, e para os economistas, em particular. 



\chapter{Cronograma}
\label{cap:cronograma}
\begin{table}[!ht]
    \centering
    \small
    \begin{tabularx}{\linewidth}{|X|c|c|c|c|c|c|c|c|c|c|c|}
    \hline
        \thead{Atividades} ~ & \thead{Mar\\2026} & \thead{Abr\\2026} & \thead{Maio\\2026} & \thead{Jun\\2026} & \thead{Jul\\2026} & \thead{Ago\\2026} & \thead{Set\\2026} & \thead{Out\\2026} & \thead{Nov\\2026} & \thead{Dez\\2026} \\ 
        \hline
        Defesa do projeto & ~ & X & X & ~ & ~ & ~ & ~ & ~ & ~ &~ \\ 
        \hline
        Trabalho bibliográfico & X & X & ~ & ~ & ~ & ~ & ~ & ~ & ~ & ~\\ 
        \hline
        Análise bibliográfica e redação capítulo 1 & ~ & X & X & ~ & ~ & ~ & ~ & ~ & ~ & ~\\ 
        \hline
        Reuniões com orientador & X & X & X & X & ~ & X & X & X & ~ & ~\\
         \hline
        Redação capítulo 2 & ~ & ~ & X & X & ~ & ~ & ~ & ~ & ~ & ~ \\ 
        \hline
        Redação capítulos 3 e 4 & ~ & ~ & ~ & X & X & X & X & ~ & ~ & ~ \\ 
        \hline
%        Redação capítulo 4 & ~ & ~ & ~ & ~ & X & X & X & ~ & ~ & ~ \\ 
%        \hline
        Redação introdução e conclusão & ~ & ~ & ~ & ~ & ~ & ~ & ~ & X & X & ~ \\ 
        \hline
        Revisões finais & ~ & ~ & ~ & ~ & ~ & ~ & ~ & ~ & X & ~ \\ 
        \hline
        Defesa da dissertação & ~ & ~ & ~ & ~ & ~ & ~ & ~ & ~ & ~ & X \\
        \hline
    \end{tabularx}
\end{table}


% \begin{table}[!ht]
%     \centering
%     \begin{longtable}{|l|l|l|l|l|l|l|l|l|l|}
%     \hline
%         Março 2026 & Abril 2026 & Maio 2026 & Junho 2026 & Julho 2026 & Agosto 2026 & Setembro 2026 & Outubro 2026 & Novembro 2026 & Dezembro 2026 \\ \hline
%         Defesa do projeto & ~ & X & X & ~ & ~ & ~ & ~ & ~ & ~ & ~ \\ \hline
%         Trabalho bibliográfico & X & X & ~ & ~ & ~ & ~ & ~ & ~ & ~ & ~ \\ \hline
%         Análise bibliográfica e redação capítulo 1 & ~ & X & X & ~ & ~ & ~ & ~ & ~ & ~ & ~ \\ \hline
%         Reuniões mensais com orientador & X & X & X & X & ~ & X & X & X & ~ & ~ \\ \hline
%         Redação capítulo 2 & ~ & ~ & X & X & ~ & ~ & ~ & ~ & ~ & ~ \\ \hline
%         Redação capítulo 3 & ~ & ~ & ~ & X & X & X & ~ & ~ & ~ & ~ \\ \hline
%         Redação capítulo 4 & ~ & ~ & ~ & ~ & X & X & X & ~ & ~ & ~ \\ \hline
%         Redação introdução e conclusão & ~ & ~ & ~ & ~ & ~ & ~ & ~ & X & ~ & ~ \\ \hline
%         Revisões finais & ~ & ~ & ~ & ~ & ~ & ~ & ~ & ~ & X & ~ \\ \hline
%         Defesa da dissertação & ~ & ~ & ~ & ~ & ~ & ~ & ~ & ~ & ~ & X \\ \hline
%     \end{longtable}
% \end{table}